\documentclass[11pt,a4paper]{article}
\usepackage[margin=0.8in]{geometry}
\usepackage{enumitem}
\usepackage{titling}

\setlength{\droptitle}{-2cm}
\title{\textbf{Internship Experience Report: SIF 2026}}
\author{
    \textbf{Yamulapalli Nandu Navadeep} \\
    \textit{IISER Thiruvananthapuram} \\[0.2cm]
    \textbf{Mentor:} Dr. Swetamber P. Das \quad \textbf{Host Institution:} SRM University, AP
}
\date{July 2026}

\begin{document}

\maketitle
\vspace{-1.5cm}

\section*{Internship Experience}
During the SIF 2026 internship program at SRM University AP, I had the invaluable opportunity to work under the supervision of Dr. Swetamber P. Das. The primary focus of my project was exploring the fundamental physical constraints on quantum dynamics, specifically the Quantum Speed Limit (QSL). The research environment was rigorous and highly stimulating, bridging the gap between abstract theoretical physics and computational application. Regular discussions and mentorship helped me navigate complex quantum derivations and significantly improved my ability to independently approach advanced theoretical problems.

\section*{Key Learnings}
The internship served as a profound learning curve, deepening my understanding of quantum mechanics far beyond the standard curriculum. My key learnings included:
\begin{itemize}[noitemsep,topsep=0pt]
    \item \textbf{Advanced Quantum Formalism:} Solidifying my grasp of Hilbert spaces, projection operators, and transitioning from pure state vectors to the density matrix formalism.
    \item \textbf{Dynamic Evolution:} Understanding the von Neumann equation as the quantum analogue to Liouville's theorem, and how density matrices temporally evolve.
    \item \textbf{Uncertainty Principles:} Distinguishing the heuristic energy-time uncertainty relation from the rigorous Robertson uncertainty relation.
    \item \textbf{Computational Physics:} Developing proficiency in Python (using \texttt{numpy} and \texttt{matplotlib}) to numerically verify theoretical bounds and visualize quantum state evolution.
\end{itemize}

\section*{Project Outcomes}
The project successfully met its research objectives, resulting in a comprehensive technical report. The major outcomes include:
\begin{itemize}[noitemsep,topsep=0pt]
    \item \textbf{Rigorous Derivation:} Successfully deriving the Mandelstam-Tamm (MT) Quantum Speed Limit from first principles, establishing the absolute minimum time required for a quantum system to reach an orthogonal state.
    \item \textbf{Theoretical Generalization:} Extending the MT bound to account for the time evolution between non-orthogonal initial and final states, characterized by the Bures angle.
    \item \textbf{Computational Verification:} Explicitly applying the generalized bound to a two-level quantum system and computationally demonstrating that the speed of quantum evolution is strictly bottlenecked by the energy variance of the system.
\end{itemize}

Overall, the SIF 2026 internship was an incredibly enriching experience that honed my theoretical and computational skills, providing a strong foundation for my future research endeavors in quantum mechanics.

\end{document}
